\documentclass[12pt,a4paper]{article}

% Pacotes
\usepackage[utf8]{inputenc}
\usepackage[T1]{fontenc}
\usepackage[brazil]{babel}
\usepackage[left=2.5cm, right=2.5cm, top=2.5cm, bottom=2.5cm]{geometry}
\usepackage{setspace}
\usepackage{titlesec}
\usepackage{parskip}
\usepackage{enumitem}
\usepackage{booktabs}
\usepackage{array}
\usepackage{xcolor}
\usepackage{hyperref}
\usepackage{fancyhdr}
\usepackage{microtype}

% Cores
\definecolor{darkblue}{HTML}{1A1A2E}
\definecolor{medgray}{HTML}{555555}
\definecolor{lightgray}{HTML}{F5F5F5}

% Formatação de seções
\titleformat{\section}{\large\bfseries\color{darkblue}}{}{0em}{}[\titlerule]
\titleformat{\subsection}{\normalsize\bfseries}{}{0em}{}
\titlespacing{\section}{0pt}{16pt}{6pt}
\titlespacing{\subsection}{0pt}{10pt}{4pt}

% Cabeçalho e rodapé
\pagestyle{fancy}
\fancyhf{}
\renewcommand{\headrulewidth}{0.4pt}
\fancyhead[L]{\small\textcolor{medgray}{Integração Vertical e Horizontal --- ADS 2026/1}}
\fancyhead[R]{\small\textcolor{medgray}{Aula 03 --- Redes Industriais}}
\fancyfoot[C]{\small\textcolor{medgray}{\thepage}}

% Espaçamento entre linhas
\setstretch{1.4}

% Hiperlinks
\hypersetup{
    colorlinks=true,
    linkcolor=darkblue,
    urlcolor=darkblue,
}

% ─────────────────────────────────────────────────────────────────────────────
\begin{document}

% Título
\begin{center}
    {\Large\bfseries\color{darkblue} Plano Orçamentário e Análise de Viabilidade}\\[4pt]
    {\large\bfseries Rede Industrial --- Indústria de Tijolos Ecológicos}\\[6pt]
    {\normalsize\textcolor{medgray}{Disciplina: Integração Vertical e Horizontal \quad|\quad ADS 2026/1}}\\[2pt]
    {\normalsize\textcolor{medgray}{Prof. Deivison Takatu}}
\end{center}

\vspace{4pt}
\hrule height 1.5pt
\vspace{14pt}

% ─────────────────────────────────────────────────────────────────────────────
\section{Contexto}

Uma indústria de produção de tijolos ecológicos irá implementar uma rede industrial para
integrar seus equipamentos de chão de fábrica com os sistemas de supervisão e gestão
corporativa. A rede deverá conectar sensores de campo, controladores lógicos programáveis
(CLPs), switches industriais, sistema supervisório (SCADA) e a rede corporativa, garantindo
comunicação determinística, robusta e compatível com futuras expansões.

Este documento apresenta o plano orçamentário com o levantamento técnico dos equipamentos
necessários, a pesquisa de fornecedores com cotações, prazos e especificações, e a
análise de viabilidade técnica e econômica da solução proposta.

% ─────────────────────────────────────────────────────────────────────────────
\section{Atividade 1 --- Plano Orçamentário e Levantamento Técnico}

\subsection{Justificativa Técnica}

Redes industriais diferem fundamentalmente das redes corporativas convencionais. Em
ambientes de manufatura, como a produção de tijolos ecológicos, a comunicação entre
dispositivos precisa ser \textbf{determinística} --- ou seja, garantir que os dados
cheguem dentro de prazos previsíveis, medidos em milissegundos. Qualquer falha de
comunicação pode interromper a linha de produção, comprometer a qualidade do produto
ou gerar riscos operacionais.

A infraestrutura proposta é baseada nos fundamentos de redes industriais estudados na
Aula 03: dispositivos de campo (sensores e atuadores), CLPs para processamento de lógica
de controle, switches industriais para interligação dos equipamentos, sistema supervisório
SCADA para monitoramento e um gateway para integração com a rede corporativa e futuros
sistemas MES/ERP.

\subsection{Levantamento dos Equipamentos}

A seguir, apresenta-se o levantamento técnico completo dos equipamentos necessários
para compor a infraestrutura de comunicação e automação da indústria.

\subsubsection*{1. Controlador Lógico Programável (CLP)}

O CLP é o núcleo do controle de processo. Ele recebe os sinais dos sensores de campo,
executa a lógica de controle programada e envia comandos aos atuadores (motores,
válvulas, prensas). Para a indústria de tijolos ecológicos, o CLP precisa suportar
entradas digitais e analógicas, comunicação via Profibus DP ou Modbus RTU/TCP e
ter capacidade de expansão por módulos.

\textbf{Equipamento selecionado:} Siemens SIMATIC S7-1200 (CPU 1214C DC/DC/DC)\\
\textbf{Quantidade:} 2 unidades (uma principal + uma redundante)\\
\textbf{Especificações técnicas:}
\begin{itemize}[leftmargin=1.5cm, itemsep=1pt]
    \item 14 entradas digitais / 10 saídas digitais integradas
    \item 2 entradas analógicas integradas
    \item Comunicação: PROFINET, Modbus TCP/RTU
    \item Tensão de alimentação: 24 VDC
    \item Temperatura de operação: 0 a 55 °C
    \item Grau de proteção: IP20
\end{itemize}
\textbf{Fornecedor:} Siemens Brasil\\
\textbf{Preço unitário:} R\$ 4.800,00\\
\textbf{Prazo de entrega:} 15 dias úteis

\subsubsection*{2. Switch Industrial}

O switch industrial é responsável por interligar todos os dispositivos da rede no chão
de fábrica. Diferente de switches comerciais, ele é projetado para suportar vibração,
ruído elétrico, poeira, umidade e variações de temperatura características de ambientes
industriais. Deve suportar o protocolo PROFINET e oferecer alta disponibilidade.

\textbf{Equipamento selecionado:} Siemens SCALANCE X208 (8 portas Fast Ethernet)\\
\textbf{Quantidade:} 2 unidades\\
\textbf{Especificações técnicas:}
\begin{itemize}[leftmargin=1.5cm, itemsep=1pt]
    \item 8 portas RJ45 10/100 Mbps
    \item Suporte a PROFINET IRT (Isochronous Real-Time)
    \item Tensão de alimentação: 24 VDC
    \item Temperatura de operação: -40 a 70 °C
    \item Grau de proteção: IP30 (montagem em trilho DIN)
    \item Diagnóstico via PROFINET e SNMP
\end{itemize}
\textbf{Fornecedor:} Siemens Brasil / distribuidora autorizada\\
\textbf{Preço unitário:} R\$ 3.200,00\\
\textbf{Prazo de entrega:} 10 dias úteis

\subsubsection*{3. Sensores de Campo}

Os sensores coletam as variáveis físicas do processo produtivo (temperatura da mistura,
pressão da prensa, nível de matéria-prima no silo e presença de peças na esteira).
Todos devem ser compatíveis com o protocolo de comunicação do CLP.

\textbf{Sensores selecionados:}

\begin{itemize}[leftmargin=1.5cm, itemsep=4pt]
    \item \textbf{Sensor de temperatura:} Balluff BTL (4--20 mA, faixa 0--200 °C,
    IP67) --- R\$ 420,00 un. --- Qtd.: 4
    \item \textbf{Sensor de pressão:} Balluff BSP (4--20 mA, 0--10 bar, IP67) ---
    R\$ 580,00 un. --- Qtd.: 3
    \item \textbf{Sensor de nível ultrassônico:} Pepperl+Fuchs UB500 (IO-Link,
    IP65, 0,2--5 m) --- R\$ 750,00 un. --- Qtd.: 2
    \item \textbf{Sensor indutivo de proximidade:} Pepperl+Fuchs NBB (NPN, 12--24 VDC,
    IP67) --- R\$ 180,00 un. --- Qtd.: 8
\end{itemize}

\textbf{Fornecedor:} Pepperl+Fuchs Brasil / Balluff Brasil\\
\textbf{Prazo de entrega:} 7 a 12 dias úteis

\subsubsection*{4. Cabeamento Estruturado Industrial}

O cabeamento industrial deve ser blindado (STP) para resistir às interferências
eletromagnéticas causadas por motores, inversores de frequência e cargas indutivas
presentes no chão de fábrica. A norma aplicável é a IEC 61784 para redes PROFINET.

\textbf{Materiais selecionados:}
\begin{itemize}[leftmargin=1.5cm, itemsep=4pt]
    \item \textbf{Cabo PROFINET Cat5e STP industrial:} SIEMENS IE FC TP Cable GP 2x2
    (rolo 100 m) --- R\$ 680,00 un. --- Qtd.: 3 rolos
    \item \textbf{Conectores RJ45 industriais blindados:} Siemens IE FC RJ45 Plug 180
    --- R\$ 45,00 un. --- Qtd.: 30
    \item \textbf{Canaletas e eletrocalhas industriais} --- R\$ 800,00 (conjunto)
    \item \textbf{Aterramento e proteção surto} --- R\$ 600,00 (conjunto)
\end{itemize}

\textbf{Fornecedor:} Siemens Brasil / materiais elétricos\\
\textbf{Prazo de entrega:} 5 dias úteis

\subsubsection*{5. Sistema Supervisório SCADA}

O sistema SCADA permite o monitoramento remoto de todos os processos, registro
histórico de variáveis, geração de alarmes e análise de desempenho. É o elo entre
o chão de fábrica e a gestão da planta.

\textbf{Software selecionado:} Siemens WinCC Runtime Professional (500 tags)\\
\textbf{Quantidade:} 1 licença\\
\textbf{Especificações:}
\begin{itemize}[leftmargin=1.5cm, itemsep=1pt]
    \item Capacidade: até 500 tags de processo
    \item Comunicação: OPC UA, PROFINET, Modbus TCP
    \item Histórico de dados: banco de dados SQL integrado
    \item Interface web para acesso remoto
    \item Compatível com TIA Portal V17
\end{itemize}
\textbf{Fornecedor:} Siemens Brasil\\
\textbf{Preço:} R\$ 8.500,00 (licença perpétua)\\
\textbf{Prazo:} Disponibilidade imediata (licença digital)

\subsubsection*{6. Gateway Industrial (OT/TI)}

O gateway é o dispositivo que integra a rede industrial (OT) com a rede corporativa (TI),
convertendo protocolos e garantindo que os dados de produção cheguem aos sistemas
de gestão (ERP/MES) de forma segura. É fundamental para viabilizar a integração vertical.

\textbf{Equipamento selecionado:} Moxa MGate MB3480 (Modbus para PROFINET/TCP)\\
\textbf{Quantidade:} 1 unidade\\
\textbf{Especificações:}
\begin{itemize}[leftmargin=1.5cm, itemsep=1pt]
    \item Conversão: Modbus RTU/ASCII/TCP para Ethernet TCP/IP
    \item 4 portas seriais RS-232/422/485
    \item Temperatura de operação: -40 a 75 °C
    \item Montagem em trilho DIN, IP30
    \item Configuração via web browser
\end{itemize}
\textbf{Fornecedor:} Moxa Brasil\\
\textbf{Preço:} R\$ 2.900,00\\
\textbf{Prazo de entrega:} 12 dias úteis

% ─────────────────────────────────────────────────────────────────────────────
\section{Atividade 2 --- Pesquisa de Fornecedores e Cotações}

\subsection{Resumo das Cotações}

A tabela abaixo consolida os fornecedores pesquisados, os equipamentos cotados, os
valores unitários, as quantidades e o custo total de cada item da infraestrutura.

\vspace{8pt}
\begin{center}
\small
\renewcommand{\arraystretch}{1.4}
\begin{tabular}{>{\bfseries}p{4cm} p{3.2cm} r r r}
\toprule
\textbf{Equipamento} & \textbf{Fornecedor} & \textbf{Vlr. Unit. (R\$)} & \textbf{Qtd.} & \textbf{Total (R\$)} \\
\midrule
CLP Siemens S7-1200        & Siemens Brasil         & 4.800,00 & 2  & 9.600,00  \\
Switch SCALANCE X208       & Siemens Brasil         & 3.200,00 & 2  & 6.400,00  \\
Sensor de temperatura      & Balluff Brasil         &   420,00 & 4  & 1.680,00  \\
Sensor de pressão          & Balluff Brasil         &   580,00 & 3  & 1.740,00  \\
Sensor de nível ultrass.   & Pepperl+Fuchs          &   750,00 & 2  & 1.500,00  \\
Sensor de proximidade      & Pepperl+Fuchs          &   180,00 & 8  & 1.440,00  \\
Cabo PROFINET STP (100 m)  & Siemens Brasil         &   680,00 & 3  & 2.040,00  \\
Conectores RJ45 industriais& Siemens Brasil         &    45,00 & 30 & 1.350,00  \\
Canaletas e eletrocalhas   & Mat. elétricos         &   800,00 & 1  &   800,00  \\
Aterramento e prot. surto  & Mat. elétricos         &   600,00 & 1  &   600,00  \\
SCADA WinCC (500 tags)     & Siemens Brasil         & 8.500,00 & 1  & 8.500,00  \\
Gateway Moxa MGate MB3480  & Moxa Brasil            & 2.900,00 & 1  & 2.900,00  \\
\midrule
\multicolumn{4}{r}{\textbf{Subtotal equipamentos}}         & \textbf{38.550,00} \\
\multicolumn{4}{r}{Instalação e comissionamento (15\%)}    &         5.782,50  \\
\multicolumn{4}{r}{Treinamento da equipe}                  &         2.000,00  \\
\multicolumn{4}{r}{Contingência (10\%)}                    &         3.855,00  \\
\midrule
\multicolumn{4}{r}{\textbf{TOTAL GERAL}}                   & \textbf{50.187,50} \\
\bottomrule
\end{tabular}
\end{center}

\vspace{8pt}

\subsection{Análise dos Fornecedores}

\textbf{Siemens Brasil} é o principal fornecedor desta solução, responsável pelos CLPs,
switches, cabeamento e software SCADA. A Siemens possui representação nacional com
suporte técnico especializado, garantia de dois anos nos equipamentos e programa de
treinamento certificado (TIA Portal). A padronização em um único ecossistema facilita
a integração entre os componentes e reduz problemas de compatibilidade.

\textbf{Balluff Brasil e Pepperl+Fuchs Brasil} são fornecedores consagrados de sensores
industriais, com linha completa certificada para ambientes com grau de proteção IP65/IP67,
compatíveis com os protocolos adotados e com estoque disponível no Brasil, o que reduz
prazos de entrega.

\textbf{Moxa Brasil} é referência global em gateways e dispositivos de comunicação
industrial, com certificações IEC e UL para ambiente industrial. O MGate MB3480 é
amplamente utilizado em integrações OT/TI e possui documentação técnica extensa.

\subsection{Adequação ao Ambiente Industrial}

Todos os equipamentos selecionados atendem às exigências do ambiente industrial da
produção de tijolos ecológicos: grau de proteção mínimo IP30 para equipamentos em
painel e IP65/IP67 para dispositivos de campo; faixa de temperatura de operação
compatível com ambientes de manufatura; resistência a vibração e ruído elétrico;
e alimentação em 24 VDC, padrão industrial.

\subsection{Escalabilidade e Compatibilidade Futura}

A solução foi desenhada com expansão em mente. O CLP Siemens S7-1200 suporta módulos
de expansão adicionais de entradas/saídas sem necessidade de substituição do hardware
principal. O switch SCALANCE X208 pode ser interligado em anel redundante (MRP) para
aumentar a disponibilidade. O SCADA WinCC pode ser migrado para versões com maior
número de tags conforme a planta cresce. O gateway Moxa garante que qualquer sistema
ERP ou MES futuro (SAP, TOTVS, etc.) possa se conectar à rede industrial via TCP/IP
sem alteração na infraestrutura de chão de fábrica.

% ─────────────────────────────────────────────────────────────────────────────
\section{Atividade 3 --- Análise de Viabilidade Técnica e Econômica}

\subsection{Viabilidade Técnica}

Do ponto de vista técnico, a solução é plenamente viável. Os equipamentos selecionados
são produtos consolidados no mercado industrial brasileiro, com histórico comprovado
em indústrias de manufatura de médio porte. A arquitetura proposta --- baseada em
PROFINET como protocolo de rede e SCADA para supervisão --- segue as melhores práticas
da Indústria 4.0 e é amplamente documentada e suportada.

A compatibilidade entre todos os componentes foi verificada: o CLP S7-1200 se comunica
nativamente com o switch SCALANCE via PROFINET, com o SCADA WinCC via OPC UA, e com
o gateway Moxa via Modbus TCP. A integração com sistemas corporativos futuros (ERP,
MES) é viabilizada pelo gateway, que expõe os dados de processo via Ethernet TCP/IP.

\subsection{Análise de Custo-Benefício}

O investimento total estimado é de \textbf{R\$ 50.187,50}. Para uma indústria de
tijolos ecológicos de médio porte, este valor é compatível com o porte da solução
e com os benefícios esperados.

Os principais ganhos operacionais identificados são:

\textbf{Redução de perdas por falha de processo:} Hoje, sem monitoramento automatizado,
falhas em sensores ou equipamentos podem passar despercebidas por horas, gerando lotes
defeituosos. Com o SCADA e os sensores integrados, alarmes são disparados em tempo
real, reduzindo perdas estimadas em 10 a 15\% do volume produzido.

\textbf{Eliminação de registros manuais:} O processo de registro de variáveis de produção
(temperatura, pressão, ciclos) deixa de ser manual e passa a ser automatizado,
reduzindo erros e liberando mão de obra para atividades de maior valor.

\textbf{Rastreabilidade e qualidade:} O histórico gerado pelo SCADA permite rastrear
os parâmetros de cada lote produzido, facilitando a identificação de causas de
não conformidades e apoiando a certificação de qualidade dos tijolos.

\textbf{Base para integração vertical:} A rede industrial implantada cria a infraestrutura
necessária para, no futuro, integrar os dados de produção ao ERP da empresa, viabilizando
planejamento integrado de produção, estoque e vendas.

\subsection{Estimativa de Retorno do Investimento (ROI)}

Considerando uma produção mensal de 50.000 tijolos ecológicos a R\$ 2,80 por unidade
(receita bruta de R\$ 140.000,00/mês), a redução estimada de perdas de 12\% representa
uma economia de R\$ 16.800,00/mês. Adicionando a redução de custos com retrabalho e
registros manuais (estimada em R\$ 2.500,00/mês), o ganho mensal total estimado é de
aproximadamente R\$ 19.300,00.

Com base nesses dados, o tempo estimado de retorno do investimento é:

\[
\text{Payback} = \frac{\text{Investimento Total}}{\text{Ganho Mensal}} = \frac{R\$\ 50.187{,}50}{R\$\ 19.300{,}00} \approx \textbf{2{,}6 \text{ meses}}
\]

Este é um payback extremamente favorável, indicando que a solução se paga em menos
de um trimestre operacional.

\subsection{Riscos e Dependências}

Os principais riscos identificados são a dependência de um único ecossistema de
fornecimento (Siemens), que pode limitar negociações futuras de preço; a necessidade
de capacitação da equipe técnica no ambiente TIA Portal, que demanda treinamento
específico; e o risco de desatualização do firmware dos equipamentos, que requer
política de atualização periódica. Esses riscos são mitigáveis por meio de contratos
de suporte e manutenção com o fornecedor, diversificação gradual de fornecedores
em expansões futuras e adoção de protocolos abertos (OPC UA, Modbus TCP) que
reduzem o lock-in tecnológico.

\subsection{Benefícios para o Negócio}

A implementação da rede industrial transforma a indústria de tijolos ecológicos em
uma operação digital integrada. Os benefícios vão além da redução de perdas: a empresa
passa a ter visibilidade total do processo produtivo em tempo real, capacidade de
tomar decisões baseadas em dados históricos e condições para crescer de forma
escalável, adicionando novos equipamentos e áreas à rede sem redesenhar a infraestrutura.
A integração com sistemas corporativos futuros posiciona a empresa para a transformação
digital, aumentando sua competitividade no mercado de construção sustentável.

% ─────────────────────────────────────────────────────────────────────────────
\vspace{20pt}
\hrule
\vspace{6pt}
{\small\textcolor{medgray}{
Disciplina: Integração Vertical e Horizontal \quad|\quad
Curso: Análise e Desenvolvimento de Sistemas \quad|\quad
ADS 2026/1 \quad|\quad
Prof. Deivison Takatu
}}

\end{document}
